% !TEX root = ../main.tex
\section{A linguagem C±}
\label{sec:S03-linguagem-cmm}

A \cmm{} (CMM) é o front-end de software do soft-processor SAPHO. O compilador \tool{cmmcomp} (\repo{yanc}; scanner \file{CMMComp.l}, parser \file{CMMComp.y}) produz \emph{assembly} simbólico (\path{.asm}) consumido pelos estágios seguintes do toolchain.

\subsection{Estrutura de um programa}

Um programa é uma sequência de elementos de topo (\code{direct}, \code{declar}, \code{funcao}), sem ordem rígida. Não há bloco \code{processor \{ ... \}}: o processador é individualizado pelas diretivas iniciais (em especial \lstinline|#PRNAME|). O ponto de entrada é \lstinline|main()|; sua ausência gera \lstinline|MSG_ERR_NO_MAIN|. A forma canônica é: diretivas no topo, depois \lstinline|main()| com laço infinito \lstinline|while (1)| modelando o hardware contínuo.

\begin{lstlisting}[style=cmm]
#PRNAME cmm_define
#NUBITS 16
#NUIOIN 1
#NUIOOU 1
#NBMANT 10
#NBEXPO 5
#NUGAIN 128
#define SCALE 3
void main() {
    int x;
    while (1) {
        x = in(0);
        x = x + SCALE;
        out(0, x);
    }
}
\end{lstlisting}

\subsection{Tipos de dados}

O scanner reconhece quatro palavras-chave de tipo, com código numérico interno (\lstinline|v_table[id].type|). Não existe tipo \enquote{ponto fixo}: o único tipo fracionário é \lstinline|float|, com precisão configurada por diretivas.

\begin{longtable}{@{}p{0.07\linewidth}p{0.14\linewidth}p{0.71\linewidth}@{}}
\toprule
\textbf{Cód.} & \textbf{Tipo} & \textbf{Descrição} \\
\midrule
\endhead
0 & \lstinline|void| & ausência de tipo \\
1 & \lstinline|int| & inteiro de \lstinline|#NUBITS| bits, complemento de dois \\
2 & \lstinline|float| & ponto flutuante (mantissa \lstinline|#NBMANT|, expoente \lstinline|#NBEXPO|) \\
3 & \lstinline|comp| & complexo (par de \emph{floats}); literal \lstinline|a+bi| (token \code{CNUM}) \\
4 & --- & parte imaginária de \lstinline|comp| (gerada automaticamente) \\
\bottomrule
\end{longtable}

Cada \lstinline|comp| ocupa dois registros (real, código 3; imaginário, código 4, criado por \lstinline|get_img_id()|). O símbolo \lstinline|i| é reservado (\lstinline|MSG_ERR_RESERVED_I|). A aritmética \lstinline|comp| suporta \lstinline|+ - * /| (módulo \file{oper.c}).

\subsection{Diretivas de hardware}

Reconhecidas como \emph{tokens} próprios; cada uma ocupa uma linha e recebe um nome ou inteiro. A regra \code{direct} chama \lstinline|dire_exec()| (\file{diretivas.c}).

\begin{longtable}{@{}p{0.16\linewidth}p{0.78\linewidth}@{}}
\toprule
\textbf{Diretiva} & \textbf{Significado} \\
\midrule
\endhead
\lstinline|#PRNAME| & nome do processador \\
\lstinline|#NUBITS| & largura da palavra da ULA \\
\lstinline|#NBMANT| & largura da mantissa (padrão 16) \\
\lstinline|#NBEXPO| & largura do expoente (padrão 6) \\
\lstinline|#NDSTAC| & profundidade da pilha de dados \\
\lstinline|#SDEPTH| & profundidade da pilha de sub-rotinas \\
\lstinline|#NUIOIN| & número de portas de entrada (padrão 1) \\
\lstinline|#NUIOOU| & número de portas de saída (padrão 1) \\
\lstinline|#NUGAIN| & constante de divisão usada por \lstinline|norm(.)| \\
\lstinline|#FFTSIZ| & tamanho da FFT ($2^{\texttt{FFTSIZ}}$) \\
\bottomrule
\end{longtable}

As comportamentais \lstinline|#PRACA| (retorno de \emph{reset} por interrupção) e \lstinline|#TOAQUI| (endereço refletido pelo pino \code{cheguei}) são \emph{statements}.

\subsection{Entrada e saída}

Portas indexadas por literal inteiro, validadas contra \lstinline|#NUIOIN|/\lstinline|#NUIOOU| (\lstinline|MSG_ERR_NO_IN_PORT|/\lstinline|_OUT_PORT|): \lstinline|in(p)| (lê), \lstinline|fin(p)| (lê convertendo para \emph{float}), \lstinline|out(p, exp);| (escreve), \lstinline|fout(p, exp);| (escreve em \emph{float}).

\subsection{Biblioteca padrão e macros}

Funções intrínsecas mapeadas a \lstinline|expr_stdlib(OP_STD_...)|: especiais (\lstinline|norm|, \lstinline|pset|, \lstinline|abs|, \lstinline|sign|, \lstinline|copy|); não-lineares e arredondamento (\lstinline|sqrt|, \lstinline|atan|, \lstinline|sin|, \lstinline|cos|, \lstinline|tan|, \lstinline|sinh|, \lstinline|cosh|, \lstinline|tanh|, \lstinline|exp|, \lstinline|log|, \lstinline|pow|, \lstinline|floor|, \lstinline|ceil|, \lstinline|round|); complexas (\lstinline|real|, \lstinline|imag|, \lstinline|fase|, \lstinline|mod2|, \lstinline|complex|, \lstinline|conj|). A \cmm{} também oferece álgebra linear em notação de Dirac (bra-ket): produto interno \lstinline+<a|b>+ como expressão, e \emph{statements} de produto matriz-vetor, externo e \emph{shift register}.

O \lstinline|#define| (apenas \emph{object-like}) é tratado no scanner com proteção contra recursão. As transcendentais são macros de \emph{assembly} em \file{Compilers/CMMComp/Includes/}, injetadas pelo codegen: \file{float\_sqrt.asm}, \file{float\_atan.asm}, \file{float\_sin.asm}, \file{float\_tan.asm}, \file{float\_exp.asm}, \file{float\_log.asm}.

\subsection{O caminho C++}

Alternativamente, o YANC compila C/C++ para o mesmo soft-processor via \tool{cpppp} (pré-processador autônomo) e \tool{cppcomp} (subconjunto de C99 com extensões de C++); os parâmetros do alvo têm padrões fixos em \file{config.h} (palavra 32 bits, mantissa 23, expoente 8) ajustáveis por \lstinline|#pragma yanc|, e a saída é o mesmo \path{.asm}.
